\chapter{Phụ lục}

\section{Phụ lục A: Danh sách bảng dữ liệu cốt lõi}
\begin{longtable}{p{4.5cm}p{9.5cm}}
\toprule
\textbf{Bảng} & \textbf{Mục đích} \\
\midrule
\texttt{tenants} & Quản lý không gian dữ liệu theo tổ chức \\
\texttt{parts} & Danh mục phụ tùng viễn thông \\
\texttt{projects} & Quản lý dự án vòng đời thiết bị \\
\texttt{project\_equipment} & Quản lý thiết bị trong từng dự án \\
\texttt{warehouses} & Quản lý kho bãi đa điểm \\
\texttt{warehouse\_stock} & Theo dõi tồn kho theo trạng thái \\
\texttt{audit\_logs} & Lưu vết thao tác quan trọng \\
\bottomrule
\end{longtable}

\section{Phụ lục B: Kịch bản kiểm thử tích hợp giai đoạn 1}
\begin{enumerate}[leftmargin=1.2cm]
    \item Tạo phụ tùng mới trong Parts Catalog.
    \item Tạo dự án ở trạng thái Draft.
    \item Thêm thiết bị từ catalog vào dự án, ghi nhận condition/disposition.
    \item Điều chuyển thiết bị vào kho đích ở trạng thái in stock.
    \item Kiểm tra KPI tổng số thiết bị và giá trị tồn kho trên dashboard.
\end{enumerate}

\section{Phụ lục C: Quy ước đánh số hình và bảng}
Trong báo cáo này, số hiệu hình và bảng được đánh theo chương:
\begin{itemize}[leftmargin=1.2cm]
    \item Hình: Chương.Số thứ tự (ví dụ: Hình 2.1, Hình 2.2).
    \item Bảng: Chương.Số thứ tự (ví dụ: Bảng 3.1, Bảng 3.2).
\end{itemize}
